This thesis evaluates three mobile pose estimation models---MediaPipe, MoveNet, and YOLOv8-Pose---plus six family variants for home-based rehabilitation. Standard benchmarks such as COCO and MPII cover general human pose estimation and do not reflect rehabilitation-specific conditions: fixed camera views, rehabilitation-style exercises, temporally stable predictions on video, and occasional secondary-person interference. We address this gap using REHAB24-6, a rehabilitation-oriented proxy benchmark of ten healthy adult volunteers performing six physiotherapist-guided exercises with synchronized RGB and motion-capture recordings.

The evaluation covers six dimensions---body-joint accuracy, viewpoint sensitivity, multi-person behavior, normalized frame-to-frame prediction displacement, detection completeness, and CPU inference speed---plus an auxiliary cross-dataset comparison on COCO val2017 under a different evaluation protocol. REHAB24-6 results are aggregated to ten subject-level clusters and analyzed with paired permutation tests, Holm correction, and bootstrap confidence intervals. MoveNet MultiPose has the lowest cluster-level body-joint error and remains significantly ahead of both alternatives after Holm correction. It also shows the lowest prediction-displacement value and the highest CPU throughput of the three main models. MediaPipe has the lowest catastrophic-failure rate on a coach extreme-case subset with a visible therapist, and YOLOv8 detects the full 12-joint skeleton most often. Detection counts alone do not explain coach-scene robustness; the three models form three architectural profiles, not a single scalar ranking. Each model is also dominated by a distinct failure mode consistent with its architectural paradigm. Model rankings on REHAB24-6 do not fully carry over to COCO under the protocol used here.

For deployment, MoveNet MultiPose is the strongest default for CPU-based body-joint rehabilitation analysis. Other configurations are better choices when detection completeness, lower coach-scene failure rate, richer native landmarks, or maximum throughput matter more. The most important next step toward clinical use is validation on real patient data.
